\documentclass[11pt]{ctexart}
\usepackage[a4paper,margin=1in]{geometry}
\usepackage{longtable,booktabs}
\title{Geant4-Agent 回归测试报告（中文）}
\author{自动化评测流程}
\date{2026-03-06\_122915}
\begin{document}
\maketitle
\section{统一 Pass/Fail 总表}
\begin{longtable}{p{0.06\linewidth}p{0.08\linewidth}p{0.07\linewidth}p{0.08\linewidth}p{0.07\linewidth}p{0.07\linewidth}p{0.07\linewidth}p{0.07\linewidth}p{0.09\linewidth}p{0.06\linewidth}p{0.21\linewidth}}
\toprule
ID & 领域 & 风格 & 参数完整性 & 期望初始 & 实际初始 & 期望最终 & 实际最终 & 缺参(初->终) & 结果 & 失败原因 \\
\midrule
MTM01 & medical & fuzzy & missing & 失败 & 失败 & 通过 & 通过 & 13->0 & 通过 & - \\
MTM02 & aerospace & precise & missing & 失败 & 失败 & 通过 & 通过 & 2->0 & 通过 & - \\
MTM03 & industrial\_ndt & precise & missing & 失败 & 失败 & 通过 & 通过 & 12->0 & 通过 & - \\
MTM04 & physics & precise & missing & 失败 & 失败 & 通过 & 通过 & 8->0 & 通过 & - \\
MTM05 & nuclear\_engineering & fuzzy & missing & 失败 & 失败 & 通过 & 失败 & 13->1 & 失败 & final\_complete\_mismatch, expected\_final\_mismatch:materials.selected\_materials.0 \\
MTM06 & security\_screening & fuzzy & missing & 失败 & 失败 & 通过 & 通过 & 10->0 & 通过 & - \\
MTM07 & semiconductor & precise & missing & 失败 & 失败 & 通过 & 失败 & 13->1 & 失败 & final\_complete\_mismatch, expected\_final\_mismatch:materials.selected\_materials.0 \\
MTM08 & education & fuzzy & missing & 失败 & 失败 & 通过 & 通过 & 11->0 & 通过 & - \\
MTM09 & environmental\_radiation & fuzzy & missing & 失败 & 失败 & 通过 & 失败 & 13->1 & 失败 & final\_complete\_mismatch, expected\_final\_mismatch:materials.selected\_materials.0 \\
MTM10 & high\_energy\_physics & precise & missing & 失败 & 失败 & 通过 & 通过 & 10->0 & 通过 & - \\
\bottomrule
\end{longtable}
\section{基准对比}
\begin{itemize}
\item pass\_rate=0.700
\item initial\_alignment\_rate=1.000
\item final\_alignment\_rate=0.700
\item delta\_pass\_rate=-0.250
\item regressions=无
\item improvements=无
\end{itemize}
\appendix
\section{附录：完整对话与内部处理链（自然语言）}
\subsection{案例 MTM01（medical）}
\textbf{对话：}
\begin{itemize}
\item 用户：????????????????????????????????
\item 系统：您想构建哪种几何结构呢？比如是一个简单的盒子，还是一个圆柱体？另外，这个结构的X方向尺寸需要设定为多少？
\item 用户：???? 300 mm x 300 mm x 300 mm ??????? G4\_WATER?source.type=beam?particle=gamma?energy=6 MeV?position=(0,0,-400) mm?direction=+z?physics=QBBC?output=root?
\item 系统：配置已完成！您设定的几何结构是一个300mm x 300mm x 300mm的盒子，材料为G4\_WATER，并配置了能量6 MeV、沿+Z方向入射的伽马射线束源，物理列表为QBBC，输出格式为ROOT文件。
\end{itemize}
\textbf{内部链路（自然语言）：}
\begin{itemize}
\item 第 1 轮：动作=ask\_clarification，阶段=geometry，闭环状态=False。
\item 缺失字段：geometry.structure、geometry.params.module\_x、geometry.params.module\_y、geometry.params.module\_z、materials.volume\_material\_map、source.type、source.particle、source.energy、source.position、source.direction、physics.physics\_list、output.format、output.path。
\item LLM 使用=True；回退=无；后端=llm\_semantic\_frame+runtime\_semantic。
\item 第 2 轮：动作=finalize，阶段=materials，闭环状态=True。
\item 缺失字段：无。
\item LLM 使用=True；回退=无；后端=llm\_slot\_frame+runtime\_semantic。
\end{itemize}
\subsection{案例 MTM02（aerospace）}
\textbf{对话：}
\begin{itemize}
\item 用户：I need a first-pass proton irradiation setup for a silicon device; keep it simple for now.
\item 系统：Here's your current setup summary:

**Geometry:** Updated with a simple structure and module dimensions.
**Materials:** Silicon assigned to your device volume.
**Source:** Proton beam configured with energy and position.
**Physics:** Appropriate physics list selected.

**To finish, please specify:**
- Output format (e.g., ROOT, CSV)
- Output file path

Once provided, I can generate your configuration.
\item 用户：Use a silicon box 50 mm x 50 mm x 5 mm, beam source proton 100 MeV from (0,0,-100) mm along +z, physics FTFP\_BERT, output json.
\item 系统：Your new geometry dimensions would overwrite existing values. Please confirm this change:

*   Module X: 1000.0 mm → 50.0 mm
*   Module Y: 1000.0 mm → 50.0 mm
\item 用户：confirm
\item 系统：Your new geometry dimensions would overwrite the existing ones. Please confirm this change:

*   Module X: 1000.0 mm → 50.0 mm
*   Module Y: 1000.0 mm → 50.0 mm
\end{itemize}
\textbf{内部链路（自然语言）：}
\begin{itemize}
\item 第 1 轮：动作=summarize\_progress，阶段=materials，闭环状态=False。
\item 缺失字段：output.format、output.path。
\item LLM 使用=True；回退=无；后端=llm\_slot\_frame+runtime\_semantic。
\item 第 2 轮：动作=confirm\_overwrite，阶段=source，闭环状态=False。
\item 缺失字段：无。
\item LLM 使用=True；回退=无；后端=llm\_slot\_frame+runtime\_semantic。
\item 第 3 轮：动作=confirm\_overwrite，阶段=physics，闭环状态=True。
\item 缺失字段：无。
\item LLM 使用=True；回退=无；后端=llm\_semantic\_frame+runtime\_semantic+confirmed\_pending\_overwrite。
\end{itemize}
\subsection{案例 MTM03（industrial\_ndt）}
\textbf{对话：}
\begin{itemize}
\item 用户：????????????????????? gamma?
\item 系统：目前配置进度如下：

**几何**：待补充结构类型与尺寸参数。  
**材料**：待指定体积与材料的绑定关系。  
**源**：已确认粒子类型为 gamma，并记录了源的选择来源和依据；还需补充源类型和能量参数。  
**物理**：待选择物理列表。  
**输出**：待设置输出格式和保存路径。

请继续补充未完成的参数，如果有需要覆盖已确认项，请明确告诉我。
\item 用户：????? 40 mm??? 80 mm??? G4\_Fe??? gamma 2 MeV??? (0,0,-200) mm??? +z?physics=Shielding?output=csv?
\item 系统：配置已完成。
\end{itemize}
\textbf{内部链路（自然语言）：}
\begin{itemize}
\item 第 1 轮：动作=summarize\_progress，阶段=geometry，闭环状态=False。
\item 缺失字段：geometry.structure、geometry.params.module\_x、geometry.params.module\_y、geometry.params.module\_z、materials.volume\_material\_map、source.type、source.energy、source.position、source.direction、physics.physics\_list、output.format、output.path。
\item LLM 使用=True；回退=无；后端=llm\_slot\_frame+runtime\_semantic。
\item 第 2 轮：动作=finalize，阶段=materials，闭环状态=True。
\item 缺失字段：无。
\item LLM 使用=True；回退=无；后端=llm\_slot\_frame+runtime\_semantic。
\end{itemize}
\subsection{案例 MTM04（physics）}
\textbf{对话：}
\begin{itemize}
\item 用户：Gamma point source 1 MeV at (0,0,-100) toward +z; please build the rest.
\item 系统：Here's your current setup:

**Source:** Configured as a gamma point source with 1 MeV energy, positioned at (0,0,-100) and directed along +z.

**What's still needed:**
*   **Geometry:** The structure type and the module size (x, y, z dimensions).
*   **Materials:** Which material to assign to your geometry volume.
*   **Physics:** Selection of a physics list.
*   **Output:** The format and file path for your results.
\item 用户：Set copper box 1 m x 1 m x 1 m, physics FTFP\_BERT, output hdf5.
\item 系统：Your configuration is complete.

**Source:** Gamma point source, 1 MeV, at (0,0,-100) directed along +z.
**Geometry:** Copper box, 1 m x 1 m x 1 m.
**Physics:** FTFP\_BERT physics list.
**Output:** HDF5 format.

All required fields are set. You can proceed with the simulation.
\end{itemize}
\textbf{内部链路（自然语言）：}
\begin{itemize}
\item 第 1 轮：动作=summarize\_progress，阶段=geometry，闭环状态=False。
\item 缺失字段：geometry.structure、geometry.params.module\_x、geometry.params.module\_y、geometry.params.module\_z、materials.volume\_material\_map、physics.physics\_list、output.format、output.path。
\item LLM 使用=True；回退=无；后端=llm\_slot\_frame+runtime\_semantic。
\item 第 2 轮：动作=finalize，阶段=materials，闭环状态=True。
\item 缺失字段：无。
\item LLM 使用=True；回退=无；后端=llm\_slot\_frame+runtime\_semantic。
\end{itemize}
\subsection{案例 MTM05（nuclear\_engineering）}
\textbf{对话：}
\begin{itemize}
\item 用户：?????????????????????
\item 系统：你希望构建什么形状的几何结构呢？另外，这个结构在 X 方向的具体尺寸是多少？
\item 用户：??????? 200 mm x 200 mm x 200 mm?source.type=point?particle=neutron?energy=14 MeV?position=(0,0,-50) mm?direction=+z?physics=QGSP\_BERT\_HP?output=xml?
\item 系统：好的，当前配置进度如下：

*   **几何**：已根据您的要求，将结构类型和X、Y方向尺寸更新为“模块”和200 mm。
*   **材料**：还需要您指定体积与材料的绑定关系。
*   **源**：已更新为点源，粒子类型为中子，能量为14 MeV。
*   **物理**：已采用您指定的 QGSP\_BERT\_HP 物理列表。
*   **输出**：已设置为生成XML格式文件。

接下来，请告诉我如何为几何体积分配材料。
\end{itemize}
\textbf{内部链路（自然语言）：}
\begin{itemize}
\item 第 1 轮：动作=ask\_clarification，阶段=geometry，闭环状态=False。
\item 缺失字段：geometry.structure、geometry.params.module\_x、geometry.params.module\_y、geometry.params.module\_z、materials.volume\_material\_map、source.type、source.particle、source.energy、source.position、source.direction、physics.physics\_list、output.format、output.path。
\item LLM 使用=True；回退=无；后端=llm\_semantic\_frame+runtime\_semantic。
\item 第 2 轮：动作=summarize\_progress，阶段=materials，闭环状态=False。
\item 缺失字段：materials.volume\_material\_map。
\item LLM 使用=True；回退=无；后端=llm\_slot\_frame+runtime\_semantic。
\end{itemize}
\subsection{案例 MTM06（security\_screening）}
\textbf{对话：}
\begin{itemize}
\item 用户：Need a quick cargo-screening toy setup with forward photons; details can follow.
\item 系统：Here's where we stand:

**Source:** Updated with your photon beam pointing forward. Still need the beam's energy and starting position.

**Geometry:** Need the overall structure and the module's X dimension.

**Materials:** Need to assign materials to the volumes.

**Physics \& Output:** Still need the physics list, output format, and file path.
\item 用户：Use steel cylinder target radius 60 mm half-length 120 mm, beam gamma 3 MeV from (0,0,-500) mm to +z, physics Shielding, output root.
\item 系统：An existing confirmed value would be overwritten. Please confirm this change: source type from "point" to "beam".
\item 用户：confirm
\item 系统：You previously confirmed a point source. Changing it to a beam source will overwrite that. Please confirm: change source type from "point" to "beam"?
\end{itemize}
\textbf{内部链路（自然语言）：}
\begin{itemize}
\item 第 1 轮：动作=summarize\_progress，阶段=geometry，闭环状态=False。
\item 缺失字段：geometry.structure、geometry.params.module\_x、geometry.params.module\_y、geometry.params.module\_z、materials.volume\_material\_map、source.energy、source.position、physics.physics\_list、output.format、output.path。
\item LLM 使用=True；回退=无；后端=llm\_slot\_frame+runtime\_semantic。
\item 第 2 轮：动作=confirm\_overwrite，阶段=materials，闭环状态=False。
\item 缺失字段：无。
\item LLM 使用=True；回退=无；后端=llm\_slot\_frame+runtime\_semantic。
\item 第 3 轮：动作=confirm\_overwrite，阶段=source，闭环状态=True。
\item 缺失字段：无。
\item LLM 使用=True；回退=无；后端=llm\_semantic\_frame+runtime\_semantic+confirmed\_pending\_overwrite。
\end{itemize}
\subsection{案例 MTM07（semiconductor）}
\textbf{对话：}
\begin{itemize}
\item 用户：?????????????????????
\item 系统：你希望构建什么形状的几何结构呢？比如一个盒子，还是一个圆柱体？另外，这个模块在X方向上的尺寸是多少？
\item 用户：?? 20 mm x 20 mm x 1 mm ?????isotropic ??? 2 MeV ? (0,0,0)?physics=QGSP\_BERT\_HP?output=json?
\item 系统：目前配置进度如下：

- **几何**：已设置结构类型和模块尺寸（X、Y方向）。
- **材料**：还需补充体积与材料的绑定关系。
- **源**：已设置类型、粒子种类和能量。
- **物理**：已选定物理列表及其依据。
- **输出**：已设定格式和保存路径。

请确认是否需要调整已同步的内容，然后我们继续补充材料部分的绑定信息。
\end{itemize}
\textbf{内部链路（自然语言）：}
\begin{itemize}
\item 第 1 轮：动作=ask\_clarification，阶段=geometry，闭环状态=False。
\item 缺失字段：geometry.structure、geometry.params.module\_x、geometry.params.module\_y、geometry.params.module\_z、materials.volume\_material\_map、source.type、source.particle、source.energy、source.position、source.direction、physics.physics\_list、output.format、output.path。
\item LLM 使用=True；回退=无；后端=llm\_semantic\_frame+runtime\_semantic。
\item 第 2 轮：动作=summarize\_progress，阶段=materials，闭环状态=False。
\item 缺失字段：materials.volume\_material\_map。
\item LLM 使用=True；回退=无；后端=llm\_slot\_frame+runtime\_semantic。
\end{itemize}
\subsection{案例 MTM08（education）}
\textbf{对话：}
\begin{itemize}
\item 用户：For classroom demo, I want gamma attenuation in copper; you can ask only essential missing parts.
\item 系统：Here's your current setup status:

**Geometry:** Still needs the structure type and module X dimension.
**Materials:** Copper is selected. Still needs to assign it to a volume.
**Source:** Gamma particle is selected. Still needs the source type and energy.
**Physics:** Physics list is configured.
**Output:** Still needs the format and file path.

To proceed, I'll need the missing geometry and source details.
\item 用户：Okay, copper box 100 mm x 100 mm x 100 mm, point gamma 1 MeV at (0,0,-200) mm, direction +z, physics FTFP\_BERT, output csv.
\item 系统：You're about to overwrite a confirmed value. Please confirm this change:

**Physics list:** Shielding → FTFP\_BERT
\item 用户：confirm
\item 系统：You're about to overwrite a confirmed value. Please confirm this change:

**Physics list:** Shielding → FTFP\_BERT
\end{itemize}
\textbf{内部链路（自然语言）：}
\begin{itemize}
\item 第 1 轮：动作=summarize\_progress，阶段=geometry，闭环状态=False。
\item 缺失字段：geometry.structure、geometry.params.module\_x、geometry.params.module\_y、geometry.params.module\_z、materials.volume\_material\_map、source.type、source.energy、source.position、source.direction、output.format、output.path。
\item LLM 使用=True；回退=无；后端=llm\_slot\_frame+runtime\_semantic。
\item 第 2 轮：动作=confirm\_overwrite，阶段=materials，闭环状态=False。
\item 缺失字段：无。
\item LLM 使用=True；回退=无；后端=llm\_slot\_frame+runtime\_semantic。
\item 第 3 轮：动作=confirm\_overwrite，阶段=source，闭环状态=True。
\item 缺失字段：无。
\item LLM 使用=True；回退=无；后端=llm\_semantic\_frame+runtime\_semantic+confirmed\_pending\_overwrite。
\end{itemize}
\subsection{案例 MTM09（environmental\_radiation）}
\textbf{对话：}
\begin{itemize}
\item 用户：???????????????????????
\item 系统：你希望构建哪种几何结构呢？比如，是一个简单的盒子，还是一个圆柱体？另外，这个结构在X方向上的尺寸是多少？
\item 用户：????? 500 mm x 500 mm x 500 mm?gamma ?? 0.5 MeV??? (0,0,-300) mm??? +z?physics=QBBC?output=root?
\item 系统：目前配置进度如下：

- **几何**：已设定结构类型和X、Y方向尺寸。
- **材料**：还需补充体积与材料的绑定关系。
- **源**：已设定类型、粒子种类和能量。
- **物理**：物理列表及相关依据已同步。
- **输出**：输出格式和路径已设定。

请确认是否需要调整已同步的项，并补充材料绑定信息。
\end{itemize}
\textbf{内部链路（自然语言）：}
\begin{itemize}
\item 第 1 轮：动作=ask\_clarification，阶段=geometry，闭环状态=False。
\item 缺失字段：geometry.structure、geometry.params.module\_x、geometry.params.module\_y、geometry.params.module\_z、materials.volume\_material\_map、source.type、source.particle、source.energy、source.position、source.direction、physics.physics\_list、output.format、output.path。
\item LLM 使用=True；回退=无；后端=llm\_semantic\_frame+runtime\_semantic。
\item 第 2 轮：动作=summarize\_progress，阶段=materials，闭环状态=False。
\item 缺失字段：materials.volume\_material\_map。
\item LLM 使用=True；回退=无；后端=llm\_slot\_frame+runtime\_semantic。
\end{itemize}
\subsection{案例 MTM10（high\_energy\_physics）}
\textbf{对话：}
\begin{itemize}
\item 用户：Set geometry as a sphere target; I have not decided radius or source details yet.
\item 系统：Here's your current setup status:

**Geometry:** Updated structure and root volume. Still need the sphere's maximum radius (child\_rmax).

**Materials:** Need to assign materials to volumes.

**Source:** Need the source type, particle type, and particle energy.

**Physics:** Need the physics list.

**Output:** Need the output format and file path.
\item 用户：Sphere target radius 80 mm, lead material, electron point source 5 MeV at (0,0,-200) mm toward +z, physics FTFP\_BERT, output hdf5.
\item 系统：Your Geant4 configuration is complete.

All required fields have been set, including the sphere's radius (80 mm), lead material, an electron point source (5 MeV), the FTFP\_BERT physics list, and HDF5 output format.
\end{itemize}
\textbf{内部链路（自然语言）：}
\begin{itemize}
\item 第 1 轮：动作=summarize\_progress，阶段=geometry，闭环状态=False。
\item 缺失字段：geometry.params.child\_rmax、materials.volume\_material\_map、source.type、source.particle、source.energy、source.position、source.direction、physics.physics\_list、output.format、output.path。
\item LLM 使用=True；回退=无；后端=llm\_slot\_frame+runtime\_semantic。
\item 第 2 轮：动作=finalize，阶段=materials，闭环状态=True。
\item 缺失字段：无。
\item LLM 使用=True；回退=无；后端=llm\_slot\_frame+runtime\_semantic。
\end{itemize}
\end{document}